\documentclass[aspectratio=169,11pt]{beamer}

\usetheme{AnnArbor}
\usepackage{fontspec}
\usepackage{amsmath,amssymb,booktabs,array,multicol}
\usepackage{listings}
\usepackage{xcolor}
\usepackage{graphicx}
\usepackage{hyperref}
\usepackage{ragged2e}

\setmainfont{DejaVu Serif}
\setsansfont{DejaVu Sans}
\setmonofont{Noto Sans Mono}

\definecolor{CodeBg}{RGB}{248,248,248}
\definecolor{TitleBlue}{RGB}{38,79,129}
\definecolor{AccentGreen}{RGB}{35,125,93}
\definecolor{SoftGray}{RGB}{80,80,80}
\definecolor{QuizOrange}{RGB}{180,95,30}

\setbeamercolor{title}{fg=TitleBlue}
\setbeamercolor{frametitle}{fg=TitleBlue}
\setbeamercolor{structure}{fg=AccentGreen}
\setbeamercolor{block title}{fg=white,bg=TitleBlue}
\setbeamercolor{block body}{bg=blue!3}
\setbeamercolor{alertblock title}{fg=white,bg=QuizOrange}
\setbeamercolor{alertblock body}{bg=orange!6}
\setbeamertemplate{navigation symbols}{}
\setbeamertemplate{footline}[frame number]
\setbeamertemplate{itemize item}{\raise0.15ex\hbox{\tiny$\blacktriangleright$}}
\setbeamertemplate{itemize subitem}{\raise0.15ex\hbox{\tiny$\bullet$}}

\lstdefinestyle{py}{
  language=Python,
  basicstyle=\ttfamily\scriptsize,
  backgroundcolor=\color{CodeBg},
  frame=single,
  framerule=0.2pt,
  rulecolor=\color{gray!35},
  breaklines=true,
  showstringspaces=false,
  keywordstyle=\color{TitleBlue}\bfseries,
  commentstyle=\color{SoftGray}\itshape,
  stringstyle=\color{AccentGreen},
  upquote=true
}

\newcommand{\key}[1]{\textbf{\textcolor{AccentGreen}{#1}}}
\newcommand{\smallnote}[1]{\vspace{0.25em}{\footnotesize\color{SoftGray}#1}}
\newcommand{\quiztitle}[1]{\textbf{\textcolor{QuizOrange}{#1}}}
\newcommand{\cmd}[1]{\texttt{\detokenize{#1}}}

\title[Introduction to Pandas]{Introduction to Pandas}
\subtitle{DataFrame, Series, Data Cleaning, Grouping, and Business Data Analysis}
\author{Duc-Minh Vu}
\institute{SLSCM Lab -- FDA -- NEU}
\date{\today}

\begin{document}

\begin{frame}
  \titlepage
\end{frame}

\begin{frame}{Learning Outcomes}
\small
By the end of this lesson, learners will be able to:
\begin{itemize}
  \item Explain the role of Pandas in data science and business analytics.
  \item Distinguish between \cmd{Series} and \cmd{DataFrame}.
  \item Create, inspect, select, filter, and sort tabular data.
  \item Read and write CSV, Excel, JSON, and text-based datasets.
  \item Detect and handle missing values, duplicate rows, and inconsistent data types.
  \item Transform columns and apply simple functions.
  \item Use \cmd{groupby()}, \cmd{agg()}, \cmd{merge()}, \cmd{concat()}, and pivot tables.
  \item Compute descriptive statistics, correlation, and simple time-series summaries.
\end{itemize}
\end{frame}

\begin{frame}{Lesson Roadmap}
\begin{columns}[T]
\column{0.48\linewidth}
\begin{enumerate}
  \item Why learn Pandas?
  \item What Pandas is
  \item Pandas in the Data Science workflow
  \item Environment setup
  \item Your first DataFrame
  \item Series and DataFrame
  \item Reading Pandas commands
\end{enumerate}
\column{0.48\linewidth}
\begin{enumerate}
  \setcounter{enumi}{7}
  \item Inspecting, selecting, filtering, and sorting
  \item Reading and writing data files
  \item Data cleaning
  \item Columns, grouping, and aggregation
  \item Merge, concat, and pivot tables
  \item Correlation, visualization, time series
  \item Practice exercises
\end{enumerate}
\end{columns}
\end{frame}

\section{What Is Pandas?}

\begin{frame}{Why Learn Pandas?}
\small
Pandas is a practical bridge between raw files and data-science models.
\begin{itemize}
  \item Real datasets are often messy, incomplete, and stored in different formats.
  \item Pandas helps organize records into rows and attributes into columns.
  \item Pandas supports inspection, cleaning, transformation, aggregation, and reporting.
  \item Pandas works well with NumPy, Matplotlib, Scikit-learn, and Statsmodels.
\end{itemize}
\begin{block}{Business analytics workflow}
Raw data $\rightarrow$ clean table $\rightarrow$ summary insight $\rightarrow$ model-ready dataset.
\end{block}
\end{frame}

\begin{frame}{What Is Pandas?}
\small
\key{Pandas} is an open-source Python library for \key{data manipulation and analysis}.
It is built on top of NumPy and provides high-level tools for structured and tabular data.
\vspace{0.4em}

Pandas is especially useful for:
\begin{itemize}
  \item reading data from CSV, Excel, JSON, and text files;
  \item cleaning, transforming, and preparing datasets;
  \item selecting, filtering, grouping, and aggregating observations;
  \item combining multiple datasets;
  \item computing descriptive statistics and quick visualizations.
\end{itemize}
\end{frame}

\begin{frame}{Pandas in the Data Science Workflow}
\small
Pandas is usually used to turn raw data into a clean, analyzable table.
\vspace{0.3em}

\begin{center}
\begin{tabular}{c c c}
\fbox{CSV / Excel / Database} & $\rightarrow$ & \fbox{Pandas DataFrame} \\[0.4em]
 & $\searrow$ & $\downarrow$ \\[0.2em]
\fbox{Raw data} & $\rightarrow$ & \fbox{Clean table} \\[0.4em]
 & $\searrow$ & $\downarrow$ \\[0.2em]
\fbox{Business question} & $\rightarrow$ & \fbox{Insight / Model} \\
\end{tabular}
\end{center}

\begin{itemize}
  \item NumPy supports efficient numerical computing.
  \item Pandas organizes data into rows, columns, and indexes.
  \item Matplotlib/Seaborn visualize results from DataFrames.
  \item Scikit-learn often receives processed data from Pandas.
\end{itemize}
\begin{block}{Key idea}
Pandas connects real-world data files with Python-based data analysis.
\end{block}
\end{frame}

\begin{frame}[fragile]{Prerequisites and Environment Setup}
\begin{block}{Recommended background}
\begin{itemize}
  \item Basic Python: variables, lists, dictionaries, loops, and functions.
  \item Basic NumPy knowledge is helpful.
  \item A working environment: Jupyter Notebook, JupyterLab, Google Colab, VS Code, or similar.
\end{itemize}
\end{block}
\begin{block}{Installation and import}
\begin{lstlisting}[style=py]
pip install pandas
\end{lstlisting}
\begin{lstlisting}[style=py]
import pandas as pd
print(pd.__version__)
\end{lstlisting}
\end{block}
\end{frame}

\begin{frame}[fragile]{Creating a DataFrame}
\small
\begin{lstlisting}[style=py]
import pandas as pd

data = {
    "Name": ["An", "Binh", "Chi"],
    "Age": [20, 21, 19],
    "GPA": [3.2, 3.6, 3.4]
}

df = pd.DataFrame(data)
print(df)
\end{lstlisting}
\begin{block}{Interpretation}
Each key in the dictionary becomes a column. Each list supplies the column values.
\end{block}
\end{frame}

\begin{frame}[fragile]{Code Output: First DataFrame}
\small
When printing \cmd{df}, Pandas displays a table with 3 rows and 3 columns.
\begin{lstlisting}[style=py]
   Name  Age  GPA
0    An   20  3.2
1  Binh   21  3.6
2   Chi   19  3.4
\end{lstlisting}
\begin{block}{What to notice}
The leftmost column \cmd{0, 1, 2} is the default index automatically created by Pandas.
\end{block}
\end{frame}

\begin{frame}{DataFrame Structure}
\small
\begin{tabular}{p{0.24\linewidth}p{0.64\linewidth}}
\toprule
\textbf{Component} & \textbf{Meaning} \\
\midrule
Rows & Records, observations, students, customers, products, or orders. \\
Columns & Variables, attributes, fields, or features. \\
Index & Row labels used for identification and selection. \\
Column labels & Names used to access variables. \\
Values & The actual data stored inside the table. \\
\bottomrule
\end{tabular}
\begin{block}{Example}
A student table may contain \cmd{StudentID}, \cmd{Name}, \cmd{Age}, \cmd{Major}, and \cmd{GPA}.
\end{block}
\end{frame}

\begin{frame}[fragile]{Series and DataFrame from the First Example}
\small
\begin{columns}[T]
\column{0.49\linewidth}
\textbf{One column is usually a Series}
\begin{lstlisting}[style=py]
print(sales["Price"])
\end{lstlisting}
\begin{lstlisting}[style=py]
0    10
1    20
2     8
Name: Price, dtype: int64
\end{lstlisting}
\column{0.49\linewidth}
\textbf{Multiple columns form a DataFrame}
\begin{lstlisting}[style=py]
print(sales[["Product", "Price"]])
\end{lstlisting}
\begin{lstlisting}[style=py]
  Product  Price
0       A     10
1       B     20
2       C      8
\end{lstlisting}
\end{columns}
\begin{alertblock}{Remember}
A labeled one-dimensional column is a \cmd{Series}; a two-dimensional table is a \cmd{DataFrame}.
\end{alertblock}
\end{frame}

\begin{frame}{How to Read Pandas Commands}
\small
Most Pandas commands follow one of these patterns:
\[
\cmd{pd.function(input, option=value)} \quad \text{or} \quad \cmd{df.method(option=value)}
\]
\begin{itemize}
  \item \cmd{pd}: standard alias for the Pandas library.
  \item \cmd{df}: a variable usually used to store a DataFrame.
  \item \cmd{function}: a package-level function, such as \cmd{pd.read_csv()}.
  \item \cmd{method}: an operation attached to a DataFrame or Series.
  \item \cmd{option=value}: named arguments controlling command behavior.
\end{itemize}
\begin{block}{Example}
\cmd{pd.read_csv("sales.csv")} reads a CSV file into a DataFrame.
\end{block}
\end{frame}

\begin{frame}{Command Pattern: Object, Attribute, Method, Function}
\small
\begin{tabular}{p{0.27\linewidth}p{0.62\linewidth}}
\toprule
\textbf{Form} & \textbf{Meaning} \\
\midrule
\cmd{pd.DataFrame(data)} & Calls a Pandas function to create a table. \\
\cmd{df.shape} & Reads an attribute describing rows and columns. \\
\cmd{df.head()} & Calls a method attached to an existing DataFrame. \\
\cmd{pd.merge(left,right)} & Calls a Pandas function to combine two tables. \\
\cmd{df["Sales"].mean()} & Selects a column, then calls a Series method. \\
\bottomrule
\end{tabular}
\vspace{0.5em}
\begin{alertblock}{Teaching note}
For every command, students should identify three parts: input, operation, and output.
\end{alertblock}
\end{frame}

\begin{frame}[fragile]{Mini Example: Explain Each Line}
\small
\begin{lstlisting}[style=py]
import pandas as pd

sales = pd.DataFrame({
    "Product": ["A", "B", "C"],
    "Quantity": [2, 3, 5],
    "Price": [10, 20, 8]
})

sales["Revenue"] = sales["Quantity"] * sales["Price"]
print(sales)
\end{lstlisting}
\begin{itemize}
  \item \cmd{pd.DataFrame(...)} creates a tabular dataset.
  \item Each dictionary key becomes a column name.
  \item \cmd{sales["Revenue"] = ...} creates a new calculated column.
  \item Multiplication is performed row by row across the selected columns.
\end{itemize}
\end{frame}

\begin{frame}[fragile]{Code Output: Revenue Example}
\small
After running the previous example, \cmd{Revenue} is created as \cmd{Quantity * Price}.
\begin{lstlisting}[style=py]
  Product  Quantity  Price  Revenue
0       A         2     10       20
1       B         3     20       60
2       C         5      8       40
\end{lstlisting}
\begin{block}{Reading the result}
Each row is one product. For example, product B has \cmd{Quantity = 3}, \cmd{Price = 20}, so \cmd{Revenue = 60}.
\end{block}
\end{frame}

\begin{frame}{Quick Check: What Is Pandas?}
\small
\quiztitle{Question 1.} Which alias is conventionally used for Pandas?
\begin{multicols}{2}
\begin{itemize}
  \item A. \cmd{pn}
  \item B. \cmd{pd}
  \item C. \cmd{ps}
  \item D. \cmd{pa}
\end{itemize}
\end{multicols}
\vspace{0.6em}
\quiztitle{Question 2.} Which Pandas structure is two-dimensional?
\begin{multicols}{2}
\begin{itemize}
  \item A. \cmd{Series}
  \item B. \cmd{tuple}
  \item C. \cmd{DataFrame}
  \item D. \cmd{ndarray}
\end{itemize}
\end{multicols}
\end{frame}

\section{DataFrame Basics}

\begin{frame}[fragile]{Inspecting a DataFrame}
\small
\begin{lstlisting}[style=py]
print(df.head())
print(df.tail())
print(df.shape)
print(df.columns)
print(df.index)
print(df.dtypes)
df.info()
print(df.describe())
\end{lstlisting}
\begin{block}{Purpose}
Inspection answers basic questions: How many rows? Which columns? Which data types? Are there missing values?
\end{block}
\end{frame}

\begin{frame}[fragile]{Example Output: Inspecting a DataFrame}
\small
For the student DataFrame above, some inspection commands return:
\begin{lstlisting}[style=py]
print(df.shape)
# (3, 3)

print(df.columns)
# Index(['Name', 'Age', 'GPA'], dtype='object')

print(df.dtypes)
# Name     object
# Age       int64
# GPA     float64
# dtype: object
\end{lstlisting}
\begin{block}{Reading the output}
\cmd{shape = (3, 3)} means the table has 3 rows and 3 columns. \cmd{object} is often used for string-like data.
\end{block}
\end{frame}

\begin{frame}{Key Commands: Inspecting DataFrames}
\small
\begin{tabular}{p{0.28\linewidth}p{0.61\linewidth}}
\toprule
\textbf{Command} & \textbf{Meaning} \\
\midrule
\cmd{df.head()} & Display the first rows. \\
\cmd{df.tail()} & Display the last rows. \\
\cmd{df.shape} & Return number of rows and columns. \\
\cmd{df.columns} & Return column labels. \\
\cmd{df.index} & Return row labels. \\
\cmd{df.dtypes} & Return the data type of each column. \\
\cmd{df.info()} & Structural summary of the DataFrame. \\
\cmd{df.describe()} & Descriptive statistics for numerical columns. \\
\bottomrule
\end{tabular}
\end{frame}

\begin{frame}[fragile]{Mini Exercise: Create and Inspect Data}
\small
Create a DataFrame named \cmd{products} with columns \cmd{Product}, \cmd{Price}, and \cmd{Stock}.
\begin{lstlisting}[style=py]
products = pd.DataFrame({
    "Product": ["Laptop", "Mouse", "Keyboard"],
    "Price": [1200, 25, 45],
    "Stock": [5, 50, 30]
})

print(products.head())
print(products.shape)
print(products.dtypes)
\end{lstlisting}
\begin{alertblock}{Check}
All columns must contain the same number of values.
\end{alertblock}
\end{frame}

\begin{frame}[fragile]{Code Output: Product DataFrame}
\small
When running the mini exercise, students can expect output like this:
\begin{lstlisting}[style=py]
    Product  Price  Stock
0    Laptop   1200      5
1     Mouse     25     50
2  Keyboard     45     30

print(products.shape)
# (3, 3)

print(products.dtypes)
# Product    object
# Price       int64
# Stock       int64
# dtype: object
\end{lstlisting}
\end{frame}

\begin{frame}{Quick Check: DataFrame Basics}
\small
\quiztitle{Question 1.} Which command creates a DataFrame?
\begin{multicols}{2}
\begin{itemize}
  \item A. \cmd{pd.DataFrame()}
  \item B. \cmd{pd.SeriesFrame()}
  \item C. \cmd{np.DataFrame()}
  \item D. \cmd{df.create()}
\end{itemize}
\end{multicols}
\vspace{0.6em}
\quiztitle{Question 2.} Which command returns the number of rows and columns?
\begin{multicols}{2}
\begin{itemize}
  \item A. \cmd{df.shape}
  \item B. \cmd{df.mean()}
  \item C. \cmd{df.merge()}
  \item D. \cmd{df.plot()}
\end{itemize}
\end{multicols}
\end{frame}

\section{Series and Data Selection}

\begin{frame}[fragile]{Creating a Series}
\small
A \key{Series} is a one-dimensional labeled array.
\begin{lstlisting}[style=py]
s = pd.Series([10, 20, 30], index=["A", "B", "C"])
print(s)
print(s.index)
print(s.values)
print(s.dtype)
\end{lstlisting}
\begin{block}{Interpretation}
A Series is like one labeled column. Pandas aligns Series by index labels during arithmetic operations.
\end{block}
\end{frame}

\begin{frame}[fragile]{Code Output: Series}
\small
\begin{lstlisting}[style=py]
A    10
B    20
C    30
dtype: int64

Index(['A', 'B', 'C'], dtype='object')
[10 20 30]
int64
\end{lstlisting}
\begin{block}{Reading the output}
The left column contains index labels; the right column contains Series values.
\end{block}
\end{frame}

\begin{frame}[fragile]{Series Alignment by Label}
\small
\begin{lstlisting}[style=py]
x = pd.Series([10, 20, 30], index=["A", "B", "C"])
y = pd.Series([1, 2, 3], index=["B", "C", "D"])

result = x + y
print(result)
\end{lstlisting}
\begin{block}{Expected idea}
Only matching labels are added directly. Unmatched labels produce missing values.
\end{block}
\end{frame}

\begin{frame}[fragile]{Code Output: Series Alignment}
\small
\begin{lstlisting}[style=py]
A     NaN
B    21.0
C    32.0
D     NaN
dtype: float64
\end{lstlisting}
\begin{block}{Why NaN?}
Labels \cmd{B} and \cmd{C} appear in both Series and are added. Label \cmd{A} appears only in \cmd{x}, and \cmd{D} appears only in \cmd{y}, so the result is \cmd{NaN}.
\end{block}
\end{frame}

\begin{frame}[fragile]{Example Data for Selection, Filtering, and Sorting}
\small
To make \cmd{loc}, \cmd{iloc}, filtering, and sorting examples clear, we use the following student table.
\begin{lstlisting}[style=py]
df = pd.DataFrame({
    "StudentID": ["S01", "S02", "S03", "S04"],
    "Name": ["An", "Binh", "Chi", "Dung"],
    "Age": [20, 21, 19, 22],
    "GPA": [3.2, 3.6, 3.4, 3.8]
})

print(df)
\end{lstlisting}
\begin{lstlisting}[style=py]
  StudentID  Name  Age  GPA
0       S01    An   20  3.2
1       S02  Binh   21  3.6
2       S03   Chi   19  3.4
3       S04  Dung   22  3.8
\end{lstlisting}
\end{frame}

\begin{frame}[fragile]{Selecting Columns}
\small
\begin{columns}[T]
\column{0.48\linewidth}
\textbf{One column returns a Series}
\begin{lstlisting}[style=py]
name = df["Name"]
print(name)
\end{lstlisting}
\column{0.48\linewidth}
\textbf{Multiple columns return a DataFrame}
\begin{lstlisting}[style=py]
subset = df[["Name", "GPA"]]
print(subset)
\end{lstlisting}
\end{columns}
\begin{block}{Syntax note}
Use one pair of brackets for one column; use a list of column names for multiple columns.
\end{block}
\end{frame}

\begin{frame}[fragile]{Code Output: Selecting Columns}
\small
\begin{columns}[T]
\column{0.48\linewidth}
\textbf{One column: Series}
\begin{lstlisting}[style=py]
0      An
1    Binh
2     Chi
3    Dung
Name: Name, dtype: object
\end{lstlisting}
\column{0.48\linewidth}
\textbf{Multiple columns: DataFrame}
\begin{lstlisting}[style=py]
   Name  GPA
0    An  3.2
1  Binh  3.6
2   Chi  3.4
3  Dung  3.8
\end{lstlisting}
\end{columns}
\end{frame}

\begin{frame}[fragile]{Selecting Rows with \cmd{loc} and \cmd{iloc}}
\small
\begin{columns}[T]
\column{0.48\linewidth}
\textbf{Label-based selection}
\begin{lstlisting}[style=py]
df = df.set_index("StudentID")
print(df.loc["S01"])
\end{lstlisting}
\column{0.48\linewidth}
\textbf{Position-based selection}
\begin{lstlisting}[style=py]
print(df.iloc[0])
print(df.iloc[0:2])
\end{lstlisting}
\end{columns}
\begin{alertblock}{Core distinction}
\cmd{loc} uses labels; \cmd{iloc} uses integer positions.
\end{alertblock}
\end{frame}

\begin{frame}[fragile]{Code Output: loc and iloc}
\small
After setting \cmd{StudentID} as the index, \cmd{loc} can select by student ID.
\begin{lstlisting}[style=py]
print(df.loc["S01"])
# Name     An
# Age      20
# GPA     3.2
# Name: S01, dtype: object

print(df.iloc[0:2])
#            Name  Age  GPA
# StudentID                
# S01          An   20  3.2
# S02        Binh   21  3.6
\end{lstlisting}
\end{frame}

\begin{frame}[fragile]{Selecting Rows and Columns Together}
\small
\begin{lstlisting}[style=py]
# label-based rows and columns
selected = df.loc[["S01", "S02"], ["Name", "GPA"]]
print(selected)

# position-based rows and columns
selected = df.iloc[0:2, [0, 2]]
print(selected)
\end{lstlisting}
\begin{block}{Reading the syntax}
The part before the comma selects rows; the part after the comma selects columns.
\end{block}
\end{frame}

\begin{frame}{Key Commands: Series and Data Selection}
\small
\begin{tabular}{p{0.30\linewidth}p{0.59\linewidth}}
\toprule
\textbf{Command} & \textbf{Meaning} \\
\midrule
\cmd{pd.Series(data)} & Create a one-dimensional labeled array. \\
\cmd{s.index} & Return Series labels. \\
\cmd{s.values} & Return underlying values. \\
\cmd{df["col"]} & Select one column as a Series. \\
\cmd{df[["c1","c2"]]} & Select multiple columns as a DataFrame. \\
\cmd{df.loc[label]} & Select by row label. \\
\cmd{df.iloc[position]} & Select by integer position. \\
\bottomrule
\end{tabular}
\end{frame}

\begin{frame}{Quick Check: Series and Selection}
\small
\quiztitle{Question 1.} A Pandas Series is:
\begin{itemize}
  \item A. a two-dimensional table
  \item B. a one-dimensional labeled array
  \item C. a database server
  \item D. a plotting package
\end{itemize}
\vspace{0.5em}
\quiztitle{Question 2.} Which selector is label-based?
\begin{multicols}{2}
\begin{itemize}
  \item A. \cmd{iloc}
  \item B. \cmd{loc}
  \item C. \cmd{shape}
  \item D. \cmd{head}
\end{itemize}
\end{multicols}
\end{frame}

\section{Filtering and Sorting Data}

\begin{frame}[fragile]{Boolean Filtering}
\small
Filter rows that satisfy a condition.
\begin{lstlisting}[style=py]
high_gpa = df[df["GPA"] >= 3.5]
print(high_gpa)
\end{lstlisting}
\begin{block}{Interpretation}
\cmd{df["GPA"] >= 3.5} creates a Boolean Series. Pandas keeps rows where the value is \cmd{True}.
\end{block}
\end{frame}

\begin{frame}[fragile]{Code Output: Filtering Data}
\small
With the condition \cmd{GPA >= 3.5}, Pandas keeps only matching rows.
\begin{lstlisting}[style=py]
high_gpa = df[df["GPA"] >= 3.5]
print(high_gpa)

#            Name  Age  GPA
# StudentID                
# S02        Binh   21  3.6
# S04        Dung   22  3.8
\end{lstlisting}
\begin{block}{Reading the output}
Students S02 and S04 have GPA greater than or equal to 3.5, so they appear in the result.
\end{block}
\end{frame}

\begin{frame}[fragile]{Filtering with Multiple Conditions}
\small
\begin{lstlisting}[style=py]
selected = df[
    (df["Age"] >= 20) &
    (df["GPA"] >= 3.5)
]
print(selected)
\end{lstlisting}
\begin{itemize}
  \item Use \cmd{&} for AND.
  \item Use \cmd{|} for OR.
  \item Use \cmd{~} for NOT.
  \item Put each condition inside parentheses.
\end{itemize}
\end{frame}

\begin{frame}[fragile]{Sorting DataFrames}
\small
\begin{lstlisting}[style=py]
# sort by one column
by_gpa = df.sort_values("GPA")

# descending order
by_gpa_desc = df.sort_values("GPA", ascending=False)

# sort by multiple columns
ordered = df.sort_values(
    ["Age", "GPA"],
    ascending=[True, False]
)
\end{lstlisting}
\begin{block}{Business example}
Sort customers by revenue descending, then by order frequency descending.
\end{block}
\end{frame}

\begin{frame}[fragile]{Code Output: Multiple Conditions and Sorting}
\small
\begin{lstlisting}[style=py]
selected = df[(df["Age"] >= 20) & (df["GPA"] >= 3.5)]
print(selected)

#            Name  Age  GPA
# StudentID                
# S02        Binh   21  3.6
# S04        Dung   22  3.8

print(df.sort_values("GPA", ascending=False))
#            Name  Age  GPA
# StudentID                
# S04        Dung   22  3.8
# S02        Binh   21  3.6
# S03         Chi   19  3.4
# S01          An   20  3.2
\end{lstlisting}
\end{frame}

\begin{frame}{Key Commands: Filtering and Sorting}
\small
\begin{tabular}{p{0.33\linewidth}p{0.56\linewidth}}
\toprule
\textbf{Command} & \textbf{Meaning} \\
\midrule
\cmd{df[df["x"] > a]} & Keep rows where column \cmd{x} exceeds \cmd{a}. \\
\cmd{(cond1) & (cond2)} & Combine conditions with AND. \\
\cmd{(cond1) | (cond2)} & Combine conditions with OR. \\
\cmd{~cond} & Negate a condition. \\
\cmd{df.sort_values("x")} & Sort rows by one column. \\
\cmd{ascending=False} & Sort in descending order. \\
\bottomrule
\end{tabular}
\end{frame}

\begin{frame}{Quick Check: Filtering and Sorting}
\small
\quiztitle{Question 1.} Which expression filters rows with \cmd{GPA >= 3.5}?
\begin{itemize}
  \item A. \cmd{df[df["GPA"] >= 3.5]}
  \item B. \cmd{df.sort_values("GPA")}
  \item C. \cmd{df.read_csv("GPA")}
  \item D. \cmd{df.merge("GPA")}
\end{itemize}
\vspace{0.5em}
\quiztitle{Question 2.} Why should parentheses be used when combining multiple filter conditions?
\end{frame}

\section{Data Input and Output}

\begin{frame}[fragile]{Reading and Writing CSV Files}
\small
\begin{columns}[T]
\column{0.48\linewidth}
\textbf{Read CSV}
\begin{lstlisting}[style=py]
df = pd.read_csv("data.csv")

df = pd.read_csv(
    "data.csv",
    sep=",",
    header=0,
    encoding="utf-8"
)
\end{lstlisting}
\column{0.48\linewidth}
\textbf{Write CSV}
\begin{lstlisting}[style=py]
df.to_csv(
    "output.csv",
    index=False
)
\end{lstlisting}
\end{columns}
\begin{alertblock}{Important}
\cmd{index=False} prevents the DataFrame index from being exported as an extra column.
\end{alertblock}
\end{frame}

\begin{frame}[fragile]{Reading and Writing Excel Files}
\small
\begin{columns}[T]
\column{0.48\linewidth}
\textbf{Read Excel}
\begin{lstlisting}[style=py]
df = pd.read_excel(
    "data.xlsx",
    sheet_name="Sheet1"
)
\end{lstlisting}
\column{0.48\linewidth}
\textbf{Write Excel}
\begin{lstlisting}[style=py]
df.to_excel(
    "output.xlsx",
    index=False
)
\end{lstlisting}
\end{columns}
\begin{block}{Teaching note}
Excel is common in business contexts, but students should still inspect types and missing values after importing.
\end{block}
\end{frame}

\begin{frame}[fragile]{JSON and Text Files}
\small
\begin{columns}[T]
\column{0.48\linewidth}
\textbf{JSON}
\begin{lstlisting}[style=py]
df = pd.read_json("data.json")

df.to_json(
    "output.json",
    orient="records"
)
\end{lstlisting}
\column{0.48\linewidth}
\textbf{Delimited text file}
\begin{lstlisting}[style=py]
df = pd.read_csv(
    "data.txt",
    sep="\t"
)
\end{lstlisting}
\end{columns}
\begin{block}{Idea}
Many structured text files can be read using \cmd{read_csv()} by changing the separator.
\end{block}
\end{frame}

\begin{frame}{Key Commands: Data I/O}
\small
\begin{tabular}{p{0.30\linewidth}p{0.59\linewidth}}
\toprule
\textbf{Command} & \textbf{Role} \\
\midrule
\cmd{pd.read_csv()} & Read CSV or delimited text files. \\
\cmd{df.to_csv()} & Write data to a CSV file. \\
\cmd{pd.read_excel()} & Read Excel files. \\
\cmd{df.to_excel()} & Write data to an Excel file. \\
\cmd{pd.read_json()} & Read JSON files. \\
\cmd{df.to_json()} & Write data to a JSON file. \\
\bottomrule
\end{tabular}
\end{frame}

\begin{frame}{Quick Check: Data I/O}
\small
\quiztitle{Question 1.} Which function reads a CSV file?
\begin{multicols}{2}
\begin{itemize}
  \item A. \cmd{pd.read_csv()}
  \item B. \cmd{pd.open_csv()}
  \item C. \cmd{pd.load_table_only()}
  \item D. \cmd{df.csv_read()}
\end{itemize}
\end{multicols}
\vspace{0.6em}
\quiztitle{Question 2.} Which argument commonly prevents the index from being exported?
\begin{multicols}{2}
\begin{itemize}
  \item A. \cmd{index=False}
  \item B. \cmd{index=True}
  \item C. \cmd{header=None}
  \item D. \cmd{drop_index=True}
\end{itemize}
\end{multicols}
\end{frame}

\section{Data Cleaning}

\begin{frame}{Why Data Cleaning Matters}
\small
Real-world datasets may contain:
\begin{itemize}
  \item missing values;
  \item duplicate rows;
  \item inconsistent data types;
  \item empty or irrelevant columns;
  \item inconsistent text, extra spaces, or inconsistent capitalization;
  \item numeric data mixed with text.
\end{itemize}
\begin{block}{Goal}
Data cleaning improves accuracy, consistency, and usability before analysis or modeling.
\end{block}
\end{frame}

\begin{frame}[fragile]{Example of Messy Data and Inspection Output}
\small
\begin{lstlisting}[style=py]
df = pd.DataFrame({
    "Name": [" An ", "Binh", "Chi", "Chi"],
    "Age": [20, None, 22, 22],
    "Price": ["10", "20", "unknown", "30"]
})

print(df)
print(df.isna().sum())
\end{lstlisting}
\begin{lstlisting}[style=py]
   Name   Age    Price
0    An   20.0       10
1  Binh    NaN       20
2   Chi   22.0  unknown
3   Chi   22.0       30

Name     0
Age      1
Price    0
dtype: int64
\end{lstlisting}
\end{frame}

\begin{frame}[fragile]{Detect Missing Values}
\small
\begin{lstlisting}[style=py]
print(df.isna())
print(df.isna().sum())

missing_by_column = df.isna().sum()
total_missing = df.isna().sum().sum()
\end{lstlisting}
\begin{block}{Interpretation}
\cmd{isna()} identifies missing cells; \cmd{sum()} counts missing values by column or across the whole DataFrame.
\end{block}
\end{frame}

\begin{frame}[fragile]{Remove or Fill Missing Values}
\small
\begin{columns}[T]
\column{0.48\linewidth}
\textbf{Remove missing values}
\begin{lstlisting}[style=py]
rows_complete = df.dropna()
cols_complete = df.dropna(axis=1)
\end{lstlisting}
\column{0.48\linewidth}
\textbf{Fill missing values}
\begin{lstlisting}[style=py]
df["Age"] = df["Age"].fillna(
    df["Age"].mean()
)

df["City"] = df["City"].fillna("Unknown")
\end{lstlisting}
\end{columns}
\begin{alertblock}{Decision question}
Should missing values be removed, imputed, or explicitly labeled? The answer depends on the analysis goal.
\end{alertblock}
\end{frame}

\begin{frame}[fragile]{Duplicates and Data Types}
\small
\begin{columns}[T]
\column{0.48\linewidth}
\textbf{Duplicate rows}
\begin{lstlisting}[style=py]
duplicate_mask = df.duplicated()
duplicate_count = duplicate_mask.sum()
df = df.drop_duplicates()
\end{lstlisting}
\column{0.48\linewidth}
\textbf{Data types}
\begin{lstlisting}[style=py]
print(df.dtypes)

df["Quantity"] = df["Quantity"].astype(int)

df["Price"] = pd.to_numeric(
    df["Price"],
    errors="coerce"
)
\end{lstlisting}
\end{columns}
\end{frame}

\begin{frame}[fragile]{String Cleaning}
\small
Text columns often need standardization.
\begin{lstlisting}[style=py]
df["Name"] = df["Name"].str.strip()
df["Name"] = df["Name"].str.lower()

df["City"] = df["City"].str.replace(
    "HN",
    "Hanoi"
)
\end{lstlisting}
\begin{block}{Common issue}
\cmd{" An "}, \cmd{"AN"}, and \cmd{"an"} may refer to the same name, but they are different strings unless cleaned.
\end{block}
\end{frame}

\begin{frame}[fragile]{Code Output: Data Cleaning}
\small
\begin{lstlisting}[style=py]
df["Age"] = df["Age"].fillna(df["Age"].mean())
df["Name"] = df["Name"].str.strip().str.lower()
df["Price"] = pd.to_numeric(df["Price"], errors="coerce")
print(df)
\end{lstlisting}
\begin{lstlisting}[style=py]
   Name    Age  Price
0    an  20.00   10.0
1  binh  21.33   20.0
2   chi  22.00    NaN
3   chi  22.00   30.0
\end{lstlisting}
\begin{block}{Reading the output}
\cmd{unknown} is converted to \cmd{NaN}; names are stripped and converted to lowercase.
\end{block}
\end{frame}

\begin{frame}{Key Commands: Data Cleaning}
\small
\begin{tabular}{p{0.32\linewidth}p{0.57\linewidth}}
\toprule
\textbf{Command} & \textbf{Meaning} \\
\midrule
\cmd{df.isna()} & Detect missing values. \\
\cmd{df.isna().sum()} & Count missing values. \\
\cmd{df.dropna()} & Remove observations with missing data. \\
\cmd{df.fillna(value)} & Replace missing values. \\
\cmd{df.duplicated()} & Detect duplicate rows. \\
\cmd{df.drop_duplicates()} & Remove duplicate rows. \\
\cmd{df.astype(type)} & Convert data type. \\
\cmd{pd.to_numeric()} & Convert values to numeric form. \\
\cmd{Series.str.*} & Apply string-processing methods. \\
\bottomrule
\end{tabular}
\end{frame}

\begin{frame}{Quick Check: Data Cleaning}
\small
\quiztitle{Question 1.} Which method removes rows containing missing values?
\begin{multicols}{2}
\begin{itemize}
  \item A. \cmd{dropna()}
  \item B. \cmd{fillna()}
  \item C. \cmd{duplicated()}
  \item D. \cmd{astype()}
\end{itemize}
\end{multicols}
\vspace{0.6em}
\quiztitle{Question 2.} Which method removes duplicate rows?
\begin{multicols}{2}
\begin{itemize}
  \item A. \cmd{drop_duplicates()}
  \item B. \cmd{dropna()}
  \item C. \cmd{sort_values()}
  \item D. \cmd{merge()}
\end{itemize}
\end{multicols}
\end{frame}

\section{Operations and Aggregation}

\begin{frame}[fragile]{Example Data for Operations and Aggregation}
\small
\begin{lstlisting}[style=py]
df = pd.DataFrame({
    "Region": ["North", "North", "South", "South", "East"],
    "Product": ["A", "B", "A", "B", "A"],
    "Quantity": [2, 3, 5, 2, 1],
    "Price": [10, 20, 8, 30, 12]
})
print(df)
\end{lstlisting}
\begin{lstlisting}[style=py]
  Region Product  Quantity  Price
0  North       A         2     10
1  North       B         3     20
2  South       A         5      8
3  South       B         2     30
4   East       A         1     12
\end{lstlisting}
\end{frame}

\begin{frame}[fragile]{Creating Calculated Columns}
\small
\begin{lstlisting}[style=py]
df["Revenue"] = df["Quantity"] * df["Price"]

total_revenue = df["Revenue"].sum()
print(df)
print(total_revenue)
\end{lstlisting}
\begin{block}{Business example}
Revenue, profit margin, discount rate, customer lifetime value, and delivery delay can often be created from existing columns.
\end{block}
\end{frame}

\begin{frame}[fragile]{Code Output: Calculated Column}
\small
\begin{lstlisting}[style=py]
  Region Product  Quantity  Price  Revenue
0  North       A         2     10       20
1  North       B         3     20       60
2  South       A         5      8       40
3  South       B         2     30       60
4   East       A         1     12       12

192
\end{lstlisting}
\begin{block}{Reading the output}
The total revenue is the sum of all row-level revenue values: 20 + 60 + 40 + 60 + 12 = 192.
\end{block}
\end{frame}

\begin{frame}[fragile]{Applying Functions}
\small
\begin{columns}[T]
\column{0.48\linewidth}
\textbf{Use \cmd{map()}}
\begin{lstlisting}[style=py]
df["Status"] = df["Score"].map(
    lambda x: "Pass" if x >= 5 else "Fail"
)
\end{lstlisting}
\column{0.48\linewidth}
\textbf{Use \cmd{apply()}}
\begin{lstlisting}[style=py]
df["Squared"] = df["Value"].apply(
    lambda x: x ** 2
)
\end{lstlisting}
\end{columns}
\begin{alertblock}{Good practice}
Prefer vectorized column operations when possible; use \cmd{apply()} when custom logic is needed.
\end{alertblock}
\end{frame}

\begin{frame}[fragile]{Normalization and Z-score Standardization}
\small
\begin{lstlisting}[style=py]
# Min-max normalization
df["Normalized"] = (
    (df["Value"] - df["Value"].min()) /
    (df["Value"].max() - df["Value"].min())
)

# Z-score standardization
df["Z"] = (
    (df["Value"] - df["Value"].mean()) /
    df["Value"].std()
)
\end{lstlisting}
\begin{block}{Interpretation}
Normalization rescales values to a fixed range; standardization expresses values in standard-deviation units.
\end{block}
\end{frame}

\begin{frame}[fragile]{Descriptive Statistics}
\small
\begin{lstlisting}[style=py]
print(df.describe())

print(df["Revenue"].mean())
print(df["Revenue"].median())
print(df["Revenue"].min())
print(df["Revenue"].max())
print(df["Revenue"].std())
\end{lstlisting}
\begin{block}{Analytical role}
Descriptive statistics help detect scale, central tendency, variation, and possible outliers.
\end{block}
\end{frame}

\begin{frame}[fragile]{Grouping with \cmd{groupby()}}
\small
\begin{lstlisting}[style=py]
region_total = df.groupby("Region")["Revenue"].sum()
print(region_total)

region_stats = (
    df.groupby("Region")["Revenue"]
      .agg(["count", "sum", "mean", "min", "max"])
)
print(region_stats)
\end{lstlisting}
\begin{block}{Split-apply-combine}
Pandas splits rows into groups, applies aggregation to each group, and combines the results.
\end{block}
\end{frame}

\begin{frame}[fragile]{Code Output: groupby and agg}
\small
\begin{lstlisting}[style=py]
Region
East      12
North     80
South    100
Name: Revenue, dtype: int64

        count  sum  mean  min  max
Region                              
East        1   12  12.0   12   12
North       2   80  40.0   20   60
South       2  100  50.0   40   60
\end{lstlisting}
\begin{block}{Reading the output}
South has the highest total revenue. The second table gives several summary statistics per region.
\end{block}
\end{frame}

\begin{frame}[fragile]{Group by Multiple Columns}
\small
\begin{lstlisting}[style=py]
summary = (
    df.groupby(["Region", "Product"])["Revenue"]
      .sum()
)
print(summary)
\end{lstlisting}
\begin{block}{Business interpretation}
This can answer questions such as: Which product contributes the most revenue in each region?
\end{block}
\end{frame}

\begin{frame}{Key Commands: Operations and Aggregation}
\small
\begin{tabular}{p{0.32\linewidth}p{0.57\linewidth}}
\toprule
\textbf{Command} & \textbf{Meaning} \\
\midrule
\cmd{df["new"] = ...} & Create or transform a column. \\
\cmd{Series.map()} & Map values to new values. \\
\cmd{Series.apply()} & Apply a function to Series elements. \\
\cmd{df.describe()} & Descriptive statistics. \\
\cmd{df.groupby()} & Group observations. \\
\cmd{.agg()} & Apply multiple aggregation functions. \\
\cmd{mean(), median(), std()} & Common numerical summaries. \\
\bottomrule
\end{tabular}
\end{frame}

\begin{frame}{Quick Check: Operations}
\small
\quiztitle{Question 1.} Which method groups observations by categories?
\begin{multicols}{2}
\begin{itemize}
  \item A. \cmd{groupby()}
  \item B. \cmd{dropna()}
  \item C. \cmd{astype()}
  \item D. \cmd{sort_index()}
\end{itemize}
\end{multicols}
\vspace{0.6em}
\quiztitle{Question 2.} Complete the idea: \cmd{Revenue = Quantity ... Price}.
\end{frame}

\section{Combining and Reshaping Data}

\begin{frame}[fragile]{Concatenating DataFrames with concat}
\small
\begin{lstlisting}[style=py]
combined = pd.concat(
    [df1, df2],
    axis=0
)
\end{lstlisting}
\begin{itemize}
  \item \cmd{axis=0}: stack rows vertically.
  \item \cmd{axis=1}: combine columns side by side.
\end{itemize}
\begin{block}{Example}
Concatenate January sales and February sales when both tables have the same columns.
\end{block}
\end{frame}

\begin{frame}[fragile]{Example Data for merge}
\small
\begin{lstlisting}[style=py]
customers = pd.DataFrame({
    "CustomerID": ["C01", "C02", "C03"],
    "Name": ["An", "Binh", "Chi"]
})

orders = pd.DataFrame({
    "OrderID": ["O01", "O02", "O03"],
    "CustomerID": ["C01", "C02", "C02"],
    "Amount": [100, 250, 300]
})
\end{lstlisting}
\begin{block}{Idea}
\cmd{CustomerID} is the common key used to combine the two tables.
\end{block}
\end{frame}

\begin{frame}[fragile]{Merging Tables with merge}
\small
\begin{lstlisting}[style=py]
result = pd.merge(
    customers,
    orders,
    on="CustomerID",
    how="inner"
)
print(result)
\end{lstlisting}
\begin{itemize}
  \item \cmd{inner}: keep matching records only.
  \item \cmd{left}: keep all records from the left table.
  \item \cmd{right}: keep all records from the right table.
  \item \cmd{outer}: keep all records from both tables.
\end{itemize}
\end{frame}

\begin{frame}[fragile]{Code Output: merge}
\small
\begin{lstlisting}[style=py]
  CustomerID  Name OrderID  Amount
0        C01    An     O01     100
1        C02  Binh     O02     250
2        C02  Binh     O03     300
\end{lstlisting}
\begin{block}{Reading the output}
Customer C02 appears twice because C02 has two orders. Customer C03 is absent in an inner merge because there is no matching order.
\end{block}
\end{frame}

\begin{frame}[fragile]{Pivot Tables}
\small
A pivot table summarizes numerical values across categories.
\begin{lstlisting}[style=py]
pivot = pd.pivot_table(
    df,
    values="Revenue",
    index="Region",
    columns="Product",
    aggfunc="sum"
)
print(pivot)
\end{lstlisting}
\begin{block}{Interpretation}
Rows are regions, columns are products, and cells contain total revenue.
\end{block}
\end{frame}

\begin{frame}[fragile]{Code Output: pivot table}
\small
For the example revenue data, a pivot table may look like this:
\begin{lstlisting}[style=py]
Product     A     B
Region             
East     12.0   NaN
North    20.0  60.0
South    40.0  60.0
\end{lstlisting}
\begin{block}{Reading the output}
The cell at row \cmd{South}, column \cmd{B} equals 60, meaning total revenue of product B in South is 60.
\end{block}
\end{frame}

\begin{frame}[fragile]{Reshaping Data: pivot and melt}
\small
\begin{columns}[T]
\column{0.48\linewidth}
\textbf{Long to wide}
\begin{lstlisting}[style=py]
wide = df.pivot(
    index="Date",
    columns="Product",
    values="Sales"
)
\end{lstlisting}
\column{0.48\linewidth}
\textbf{Wide to long}
\begin{lstlisting}[style=py]
long = pd.melt(
    wide.reset_index(),
    id_vars="Date"
)
\end{lstlisting}
\end{columns}
\begin{block}{Idea}
The data shape should match the analytical task: visualization, aggregation, modeling, or reporting.
\end{block}
\end{frame}

\begin{frame}{Key Commands: Combine and Reshape}
\small
\begin{tabular}{p{0.32\linewidth}p{0.57\linewidth}}
\toprule
\textbf{Command} & \textbf{Meaning} \\
\midrule
\cmd{pd.concat()} & Combine DataFrames along an axis. \\
\cmd{pd.merge()} & Combine DataFrames using key columns. \\
\cmd{how="inner"} & Keep matching records only. \\
\cmd{how="left"} & Keep all records from the left table. \\
\cmd{df.pivot()} & Reshape long data to wide format. \\
\cmd{pd.melt()} & Reshape wide data to long format. \\
\cmd{pd.pivot_table()} & Create a summarized table. \\
\bottomrule
\end{tabular}
\end{frame}

\begin{frame}{Quick Check: Combine and Reshape}
\small
\quiztitle{Question 1.} Which function combines DataFrames using a key column?
\begin{multicols}{2}
\begin{itemize}
  \item A. \cmd{pd.merge()}
  \item B. \cmd{pd.mean()}
  \item C. \cmd{pd.reshape()}
  \item D. \cmd{pd.filter_rows()}
\end{itemize}
\end{multicols}
\vspace{0.6em}
\quiztitle{Question 2.} What does a pivot table do?
\end{frame}

\section{Advanced Operations}

\begin{frame}[fragile]{Correlation}
\small
Correlation measures association between numerical variables.
\begin{lstlisting}[style=py]
print(df.corr(numeric_only=True))

print(
    df["Advertising"].corr(df["Sales"])
)
\end{lstlisting}
\begin{alertblock}{Important}
Correlation does not by itself imply causation.
\end{alertblock}
\end{frame}

\begin{frame}[fragile]{Code Output: Correlation and Time Series}
\small
\begin{lstlisting}[style=py]
print(df[["Revenue", "Quantity", "Price"]].corr())
\end{lstlisting}
\begin{lstlisting}[style=py]
           Revenue  Quantity     Price
Revenue   1.000000  0.620174  0.594089
Quantity  0.620174  1.000000 -0.272772
Price     0.594089 -0.272772  1.000000
\end{lstlisting}
\begin{block}{Reading the output}
The correlation matrix is symmetric. Values closer to 1 indicate a stronger positive linear relationship; however, correlation does not prove causation.
\end{block}
\end{frame}

\begin{frame}[fragile]{Quick Visualization with Pandas}
\small
Pandas plotting functions are built on top of Matplotlib.
\begin{columns}[T]
\column{0.48\linewidth}
\begin{lstlisting}[style=py]
df.plot(x="Month", y="Sales", kind="line")

df.plot(x="Product", y="Sales", kind="bar")
\end{lstlisting}
\column{0.48\linewidth}
\begin{lstlisting}[style=py]
df["Sales"].plot(kind="hist")

df.plot(x="Advertising", y="Sales", kind="scatter")
\end{lstlisting}
\end{columns}
\end{frame}

\begin{frame}[fragile]{Time-Series Pipeline}
\small
\begin{lstlisting}[style=py]
df["Date"] = pd.to_datetime(df["Date"])
df = df.set_index("Date")
df = df.sort_index()

weekly = df["Sales"].resample("W").sum()
monthly = df["Sales"].resample("ME").sum()
\end{lstlisting}
\begin{block}{Time-series idea}
After converting the date column and setting it as the index, Pandas can aggregate data by calendar periods.
\end{block}
\end{frame}

\begin{frame}[fragile]{Datetime Components and Moving Average}
\small
\begin{columns}[T]
\column{0.48\linewidth}
\textbf{Date components}
\begin{lstlisting}[style=py]
df["Year"] = df.index.year
df["Month"] = df.index.month
df["Day"] = df.index.day
\end{lstlisting}
\column{0.48\linewidth}
\textbf{Moving average}
\begin{lstlisting}[style=py]
df["MovingAvg"] = (
    df["Sales"]
      .rolling(window=3)
      .mean()
)
\end{lstlisting}
\end{columns}
\end{frame}

\begin{frame}{Key Commands: Advanced Operations}
\small
\begin{tabular}{p{0.32\linewidth}p{0.57\linewidth}}
\toprule
\textbf{Command} & \textbf{Meaning} \\
\midrule
\cmd{df.corr()} & Calculate correlation matrix. \\
\cmd{Series.corr()} & Calculate correlation between two Series. \\
\cmd{df.plot()} & Create quick visualizations. \\
\cmd{pd.to_datetime()} & Convert values to datetime. \\
\cmd{df.set_index()} & Set a column as the index. \\
\cmd{df.resample()} & Aggregate time-series data by time intervals. \\
\cmd{Series.rolling()} & Calculate rolling-window statistics. \\
\bottomrule
\end{tabular}
\end{frame}

\begin{frame}{Quick Check: Advanced Operations}
\small
\quiztitle{Question 1.} Which function converts a column to datetime?
\begin{multicols}{2}
\begin{itemize}
  \item A. \cmd{pd.to_datetime()}
  \item B. \cmd{pd.date_convert_only()}
  \item C. \cmd{df.datetime()}
  \item D. \cmd{pd.time_series()}
\end{itemize}
\end{multicols}
\vspace{0.5em}
\quiztitle{Question 2.} Which method calculates a correlation matrix?
\begin{multicols}{2}
\begin{itemize}
  \item A. \cmd{corr()}
  \item B. \cmd{merge()}
  \item C. \cmd{dropna()}
  \item D. \cmd{pivot()}
\end{itemize}
\end{multicols}
\end{frame}

\section{Summary and Practice}

\begin{frame}{Content Summary I}
\small
\begin{tabular}{p{0.27\linewidth}p{0.63\linewidth}}
\toprule
\textbf{Topic} & \textbf{Main idea} \\
\midrule
Pandas & Python library for structured data manipulation and analysis. \\
Series & One-dimensional labeled array. \\
DataFrame & Two-dimensional labeled table. \\
Inspection & Use \cmd{head()}, \cmd{shape}, \cmd{info()}, and \cmd{describe()}. \\
Selection & Use column brackets, \cmd{loc}, and \cmd{iloc}. \\
Filtering & Use Boolean conditions to keep relevant rows. \\
\bottomrule
\end{tabular}
\end{frame}

\begin{frame}{Content Summary II}
\small
\begin{tabular}{p{0.27\linewidth}p{0.63\linewidth}}
\toprule
\textbf{Topic} & \textbf{Main idea} \\
\midrule
Data I/O & Read and write CSV, Excel, JSON, and text files. \\
Cleaning & Handle missing values, duplicates, data types, and strings. \\
Operations & Create calculated columns and transform data. \\
Aggregation & Use \cmd{groupby()} and \cmd{agg()} for summary reports. \\
Combination & Use \cmd{concat()} and \cmd{merge()} for multiple tables. \\
Advanced & Correlation, visualization, and time-series operations. \\
\bottomrule
\end{tabular}
\end{frame}

\begin{frame}{Practice: Student Data}
\small
\begin{block}{Exercise 1}
Create a DataFrame with at least 10 students containing:
\begin{itemize}
  \item \cmd{StudentID}, \cmd{Name}, \cmd{Age}, \cmd{Major}, \cmd{GPA}.
\end{itemize}
Perform:
\begin{enumerate}
  \item inspection with \cmd{head()}, \cmd{shape}, \cmd{info()}, and \cmd{describe()};
  \item filtering for \cmd{GPA >= 3.5};
  \item sorting by GPA descending;
  \item selecting rows using \cmd{loc} and \cmd{iloc}.
\end{enumerate}
\end{block}
\end{frame}

\begin{frame}{Practice: Missing Data and Cleaning}
\small
\begin{block}{Exercise 2}
Add missing values to the student dataset. Then:
\begin{enumerate}
  \item count missing values;
  \item fill missing Age with mean Age;
  \item fill missing GPA with median GPA;
  \item remove rows with missing Name.
\end{enumerate}
\end{block}
\begin{block}{Exercise 3}
Create a messy dataset containing spaces in names, duplicate rows, and mixed numeric/text values. Clean it using \cmd{str.strip()}, \cmd{drop_duplicates()}, and \cmd{pd.to_numeric()}.
\end{block}
\end{frame}

\begin{frame}{Practice: Sales Analysis}
\small
\begin{block}{Exercise 4}
Create a sales dataset containing \cmd{Date}, \cmd{Region}, \cmd{Product}, \cmd{Quantity}, and \cmd{Price}.
Create:
\[
\text{Revenue} = \text{Quantity} \times \text{Price}.
\]
Then calculate:
\begin{enumerate}
  \item total revenue;
  \item revenue by Region;
  \item revenue by Product;
  \item average revenue by Region;
  \item a Region $\times$ Product pivot table.
\end{enumerate}
\end{block}
\end{frame}

\begin{frame}{Practice: Merge and Time Series}
\small
\begin{block}{Exercise 5: Merge}
Create \cmd{customers(CustomerID, Name, City)} and \cmd{orders(OrderID, CustomerID, Amount)}. Perform inner merge, left merge, and identify customers without orders.
\end{block}
\begin{block}{Exercise 6: Time series}
Create daily sales data for at least 30 days. Convert \cmd{Date} to datetime, set it as index, calculate weekly total sales, calculate a 7-day moving average, and create a line plot.
\end{block}
\end{frame}

\begin{frame}{Final Reflection}
\small
\begin{itemize}
  \item What is the difference between a Series and a DataFrame?
  \item When should we use \cmd{loc}, and when should we use \cmd{iloc}?
  \item Why should we inspect \cmd{shape}, \cmd{dtypes}, and missing values early?
  \item What is the difference between \cmd{concat()} and \cmd{merge()}?
  \item Why is \cmd{groupby()} important in business analytics?
  \item Why does correlation not automatically prove causation?
\end{itemize}
\end{frame}

\begin{frame}{Suggested Answers: Quick Checks}
\small
\begin{tabular}{p{0.42\linewidth}p{0.45\linewidth}}
\toprule
\textbf{Section} & \textbf{Answers} \\
\midrule
What is Pandas? & B; C \\
DataFrame basics & A; A \\
Series and selection & B; B \\
Data I/O & A; A \\
Data cleaning & A; A \\
Operations & A; Revenue = Quantity $\times$ Price \\
Combine and reshape & A; summarizes values across categories \\
Advanced operations & A; A \\
\bottomrule
\end{tabular}
\end{frame}

\end{document}
